\chapter{Related Work}
\section{Noise Models and Robust Learning}
Fr\'{e}nay and Verleysen organize label-noise research around the assumptions
that connect the observed and true labels~\cite{frenay2014}. That distinction
matters because a uniform random flip is not equivalent to systematic
confusion between particular classes or feature-dependent annotation errors.
The present benchmark implements only uniform alternative-label flips and
fixed cyclic transitions. Its controlled corruption mask makes detection
measurable, but does not make either mechanism a realistic model of all
annotation processes.

Li and Zhu review deep-learning approaches to noisy-label image
classification and examine noise patterns informed by real-world
data~\cite{li2024}. The current work occupies the narrower
sample-selection and label-replacement part of this landscape. It does not
adapt the training loss or architecture. This restriction allows the same
cleaning interface to be used with linear, bagged-tree, and boosted-tree
estimators, but leaves stronger robust-training approaches outside the
experimental comparison.

\section{Evidence Specific to Tabular Data}
Nazaretyan et al. survey 34 approaches and benchmark 20 methods for identifying
mislabeled tabular samples~\cite{nazaretyan2025}. Their evaluation includes
controlled corruption and a genomic dataset with known errors. The breadth of
that study is greater than the present project. This project's role is a
small, auditable comparison of inexpensive probability-based scores and the
consequences of acting on their rankings. Its detection values should not be
compared directly with theirs: datasets, noise processes, selection budgets,
and learners are different.

Yang et al. investigate whether noisy-label methods associated with computer
vision can help deep learning on electronic health records, and report
beneficial adaptations in their evaluated settings~\cite{yang2024}. Their
work motivates testing rather than assuming transfer across domains. This
project uses neither their EHR data nor their deep networks; its clinical
dataset is a small image-derived diagnostic table. It therefore cannot
validate clinical deployment or reproduce their domain-specific findings.

Eisenb{\"u}rger et al. study label noise in gradient-boosted decision trees,
including detection, relabeling, and early stopping~\cite{eisenburger2025}.
They motivate including a boosted-tree learner, but this benchmark uses
scikit-learn's histogram gradient boosting rather than implementing their
Gradients detector. No claim of equivalence to their method or their reported
performance is made. Their preprint's version history starts in September
2024; the January 2025 version is the one cited here.

\section{Confidence Is Not Confident Learning}
Northcutt, Jiang, and Chuang's confident learning estimates the relationship
between observed and latent labels using probabilistic thresholds and a
class-conditional noise model~\cite{northcutt2021}. It combines ranking with
noise-distribution estimation. In contrast, this project's confidence score
is simply one minus the out-of-fold probability of the observed label.
Its cleaning budget is supplied from the experimental noise rate. It does
not estimate the confident joint, implement the cleanlab algorithm, or establish
the theoretical guarantees of confident learning. Confident learning is an
important omitted comparator for future work.

\section{Ground Truth and Statistical Evidence}
De Vries and Thierens identify a difficulty with synthetic-noise experiments:
adding known corruptions to a real dataset does not establish that its original
labels were clean~\cite{devries2025}. Their SYNLABEL framework constructs a
ground-truth function and soft-label representations to make that distinction
explicit. This project instead retains original public labels. Consequently,
its detection metrics measure recovery of \emph{injected changes}, not discovery
of every genuinely incorrect original annotation.

Dem{\v{s}}ar emphasizes comparisons across datasets when evaluating learning
algorithms~\cite{demsar2006}. This project adopts datasets as the aggregate
inferential unit, not individual rows of its result matrix. It uses an exact
paired sign-flip test rather than reproducing Dem{\v{s}}ar's recommended
testing suite. The four datasets provide only limited inferential resolution;
overlapping random splits and multiple conditions do not create hundreds of
independent datasets.

\section{Position of This Project}
The distinguishing emphasis is traceability: each result can be connected to
the split, corruption count, detector, cleaning action, seed, dependencies,
and training revision that produced it. The comparison holds these elements
fixed while changing the score or cleaning action. It is a focused empirical
study, not a comprehensive competition against the literature. Its strongest
supported conclusions concern relative behavior within this matrix and the
failure of an intuitively appealing weighted score to beat a simpler baseline.
